\begin{table}[htbp]
\centering
\small
\caption{Resumen de optimización intermedia y rol de cada variante.}
\label{tab:ch4_intermediate_optimization_summary}
\begin{tabular}{@{}p{2.9cm}p{4.7cm}p{4.0cm}@{}}
\toprule
Variante o etapa & Resultado experimental que aporta & Uso en la comparación final \\
\midrule
Búsqueda por familia & Permite identificar rangos plausibles de hiperparámetros para métodos candidatos sin fijar conclusiones definitivas. & Evidencia intermedia para reducir el espacio de búsqueda. \\
Configuración global & Evalúa si una parametrización única puede sostener desempeño razonable entre escenarios. & Favorece reproducibilidad y evita reajuste por caso. \\
TRSS fijo & Representa una variante estable con parámetros congelados antes de la comparación pareada. & Lectura principal frente a MNE. \\
TRSS calibrado & Aproxima una selección más adaptativa, pero aún depende de reglas de calibración. & Análisis auxiliar de sensibilidad. \\
Oráculo & Estima una cota superior no desplegable al escoger la mejor variante observada. & Contexto metodológico; no se reporta como método aplicable. \\
\bottomrule
\end{tabular}
\end{table}
